\doublespacing % Do not change - required

\chapter{Introduction}
\label{ch:introduction}

%%%%%%%%%%%%%%%%%%%%%%%%%%%%%%%%%%%%%%%
% IMPORTANT
\begin{spacing}{1} %THESE FOUR
\minitoc % LINES MUST APPEAR IN
\end{spacing} % EVERY
\thesisspacing % CHAPTER
% COPY THEM IN ANY NEW CHAPTER
%%%%%%%%%%%%%%%%%%%%%%%%%%%%%%%%%%%%%%%

\section{Motivation}

Use this chapter to introduce the problem addressed by the thesis as a whole, rather than by any single paper. The opening should explain the practical or scientific setting, why the problem matters, and what gap remains despite existing work.

\section{Research questions}

State the thesis-level research questions here. For a thesis built from publications, these questions should make clear why the four paper chapters belong in one coherent doctoral project.

\section{Thesis contributions}

Summarise the main contributions of the thesis. This section should synthesise across the publications instead of repeating each paper abstract.

\section{Structure of the thesis}

Briefly explain how the thesis is organised. Chapter~\ref{ch:literature-review} should review the shared research landscape and position the thesis before the publication-based chapters begin. For example, Chapter~\ref{ch:paper1} may introduce the first technical contribution, Chapter~\ref{ch:paper2} may extend the method or evaluation, Chapter~\ref{ch:paper3} may address a complementary challenge, and Chapter~\ref{ch:paper4} may provide the final study or system contribution. Chapter~\ref{ch:synthesis} should then connect the findings across the four publications.
